% Compilen desde esta carpeta o desde la raíz del repositorio.
\makeatletter
\def\input@path{{../../preamble/}{book/preamble/}}
\makeatother
\input{preamble.tex}

\usepackage[spanish]{babel}
\usepackage{lmodern}
\usepackage{titlesec}
\geometry{a4paper,margin=1.6cm}
\linespread{1.10}

\definecolor{HeadingColor}{HTML}{243B53}
\definecolor{InstructionColor}{HTML}{704214}
\definecolor{ConceptColor}{HTML}{205C50}
\definecolor{LinkColor}{HTML}{0047AB}
\hypersetup{allcolors=LinkColor}
\newcommand{\guidelink}[2]{\href{#1}{\textbf{\underline{#2}}}}
\newcommand{\instruction}[1]{\textcolor{InstructionColor}{\textbf{#1}}}
\newcommand{\concept}[1]{\textcolor{ConceptColor}{\textbf{#1}}}
\newcommand{\checkbox}{%
    \begingroup
    \setlength{\fboxrule}{0.4pt}%
    \setlength{\fboxsep}{2pt}%
    \fbox{\rule{0pt}{0.8ex}\rule{0.8ex}{0pt}}%
    \endgroup}

\setlength{\parskip}{4pt}
\setlength{\tabcolsep}{5pt}
\renewcommand{\arraystretch}{1.15}
\setlist{leftmargin=0.65cm,topsep=4pt,partopsep=0pt,parsep=0pt,itemsep=3pt}
\titleformat{\section}{\large\bfseries\color{HeadingColor}}{}{0pt}{}
\titlespacing*{\section}{0pt}{10pt}{3pt}
\newcolumntype{Y}{>{\raggedright\arraybackslash}X}

\hypersetup{pdftitle={C9--C10. Reflexionen y sustenten el cierre del proyecto},pdfsubject={Emprendimiento e Innovación, Grado 10, Periodo 2}}

\begin{document}

\begin{center}
{\Large\bfseries\color{HeadingColor} C9--C10. Reflexionen y sustenten el cierre del proyecto\par}
\vspace{3pt}
Emprendimiento e Innovación \enspace|\enspace Grado 10 \enspace|\enspace Periodo 2\\
Guía 5 de 5
\end{center}

\section*{Objetivo y punto de partida}
\concept{C9.} Coopera en el equipo aportando ideas, cumpliendo compromisos y resolviendo desacuerdos de manera constructiva.\\
\concept{C10.} Explica el valor diferencial de la solución frente a la necesidad inicialmente identificada.

Revisen el trabajo desarrollado hasta C8: D1 y los H de P1; coevaluación y selección, representación, versiones, decisiones, seguimiento, pruebas y ajustes de P2. El propósito es hacer un balance verificable del proyecto y de la colaboración, explicar su valor y analizar retroalimentación final.

\section*{Documento acumulativo e integridad de la evidencia}
Continúen el \textbf{documento acumulativo en LaTeX} en \guidelink{https://www.overleaf.com/}{Overleaf} o su editor habitual; entreguen fuente y PDF actualizado. Distingan \textbf{P1} y \textbf{P2}. Enlacen evidencias, decisiones, versiones y resultados por códigos o referencias internas, con fuentes y fechas reales.

Distingan \concept{observaciones y comentarios}, \concept{interpretaciones}, \concept{supuestos} e \concept{ideas}; identifiquen simulaciones y límites. \instruction{Usen IA solo para código LaTeX, formato y organización de contenido propio.} No inventen evidencias.

Actualicen la \textbf{bitácora} por sesión: fecha, aportes, evidencias, compromisos, responsables, plazos, estados y acuerdos. Las guías admiten proyectos físicos, digitales, comunicativos, de servicios o procedimentales. Adapten medios y extensión a la evidencia necesaria; los formatos sugeridos son opciones.

\section*{C9 de P2. Reconstruyan y valoren la colaboración}
\begin{enumerate}[series=pasos]
\item \textbf{Revisen la bitácora y contrasten sus registros.} Identifiquen aportes de cada integrante, compromisos, estados y resultados. Relaciónenlos con alternativas, decisiones, construcción, pruebas o ajustes. Si faltan registros, declaren la limitación; cualquier reconstrucción posterior debe identificarse como tal y citar su respaldo, sin inventar entradas pasadas.
\item \textbf{Sinteticen cómo trabajaron juntos.} Expliquen quién aportó qué y qué ocurrió con los compromisos: cumplidos, reprogramados o pendientes, con sus razones. Describan dificultades, acuerdos y cómo resolvieron desacuerdos de manera constructiva. Si no hubo un desacuerdo registrado, expliquen cómo escucharon razones para acordar una decisión real.
\item \textbf{Reflexionen con evidencia.} Relacionen los acuerdos y dificultades con sus efectos en el trabajo realizado. Expliquen qué aprendieron sobre cooperación, organización y comunicación, qué práctica les ayudó y qué podrían mejorar en otro proyecto. Distingan el hecho registrado de la interpretación del equipo y recojan las perspectivas de sus integrantes.
\end{enumerate}

\section*{Entregables de C9}
\begin{itemize}
\item \textbf{Síntesis verificable de colaboración en C9.} Presenten los aportes, compromisos y sus estados, acuerdos, dificultades y aprendizajes por integrante y como equipo, con códigos o referencias a la bitácora y a los productos. Expliquen cómo la colaboración influyó en los resultados y recojan las perspectivas de sus integrantes.
\item \textbf{Bitácora completa y declaración de vacíos.} Conserven los registros que respaldan la síntesis e identifiquen las limitaciones o reconstrucciones posteriores con sus fuentes; una valoración general como «trabajamos bien» sin respaldo no evidencia el criterio.
\end{itemize}

\newpage
\section*{C10 de P2. Expliquen el valor y analicen retroalimentación final}
\begin{enumerate}[resume=pasos]
\item \textbf{Construyan una explicación comparativa.} Retomen usuario y necesidad inicial de D1 y sus H. Comparen la solución disponible al cierre de C8 con la forma previa de atender la necesidad o con alternativas identificadas en C1--C2. Expliquen qué aporta de diferente, para quién y mediante qué característica, función o interacción concreta.
\item \textbf{Sustenten el alcance del valor.} Relacionen cada afirmación con rasgos demostrables del prototipo, resultados de pruebas y ajustes hasta C8. Identifiquen las versiones: distingan lo observado en la versión probada del beneficio esperado de un ajuste posterior. Expliquen limitaciones y condiciones de viabilidad; si algo no se comprobó, indíquenlo.
\item \textbf{Presenten el balance y reciban retroalimentación.} Compartan la explicación y las evidencias pertinentes con pares o docentes mediante una demostración, conversación, recorrido comentado u otro formato adecuado. Registren fecha, quién retroalimenta o su perfil, fortalezas, aspectos por mejorar y posibles pasos futuros, con referencia al trabajo presentado. Esta devolución final no reemplaza la prueba de C7 ni demuestra por sí sola el funcionamiento de una versión.
\item \textbf{Analicen y cierren.} Separen los comentarios recibidos de la interpretación del equipo. Expliquen qué aceptan, cuestionan o dejan pendiente y por qué, contrastándolo con sus evidencias. Cierren distinguiendo \textbf{lo realizado, lo aprendido y lo pendiente}. Los posibles pasos futuros son recomendaciones razonadas; \textbf{no constituyen una nueva fase obligatoria de construcción} ni deben presentarse como mejoras ya implementadas.
\end{enumerate}

\section*{Entregables de C10}
\begin{itemize}
\item \textbf{Explicación del valor diferencial en C10.} Incluyan necesidad inicial, usuario y referente de comparación; sustenten qué aporta la solución mediante evidencias del prototipo, pruebas y ajustes, con versiones y límites. Distingan resultados observados de beneficios esperados.
\item \textbf{Registro y análisis de retroalimentación final.} Identifiquen fecha, fuentes y trabajo presentado; separen comentarios e interpretaciones y expliquen qué aceptan, cuestionan o dejan pendiente y por qué. Incluyan fortalezas, aspectos por mejorar y posibles pasos futuros, contrastados con sus evidencias.
\item \textbf{Balance de cierre y soporte de presentación.} Distingan lo realizado, lo aprendido y lo pendiente, sin presentar recomendaciones futuras como mejoras implementadas. Conserven el soporte usado para presentar el trabajo o una descripción que permita reconstruir qué mostraron.
\end{itemize}

\section*{Resultado de cierre}
Entreguen el documento acumulativo en LaTeX y su PDF actualizado con las secciones C1--C10, la bitácora y referencias legibles a los soportes y versiones disponibles. El cierre sintetiza el trabajo hasta C8 y la reflexión de C9--C10; no requiere preparar una feria ni construir una versión adicional.

\section*{Pregunta de reflexión}
¿Cómo influyeron sus formas de colaborar en el valor que hoy pueden sustentar de la solución, a partir de las decisiones, resultados y dificultades documentados durante el proyecto?

Incorporen su respuesta al documento acumulativo y susténtenla en la experiencia y las evidencias del equipo.

\end{document}
